%=======================================================================
%  GuardianLink -- Capstone Video Run Sheet
%  Compile:  pdflatex GuardianLink-VideoRunSheet.tex  (run twice for TOC)
%=======================================================================
\documentclass[11pt,letterpaper]{article}

\usepackage[T1]{fontenc}
\usepackage[utf8]{inputenc}
\usepackage[margin=1in]{geometry}
\usepackage{lmodern}
\usepackage{microtype}
\usepackage{amssymb}
\usepackage{booktabs}
\usepackage{tabularx}
\usepackage{longtable}
\usepackage{array}
\usepackage{enumitem}
\usepackage{fancyvrb}
\usepackage{float}
\usepackage{fancyhdr}
\usepackage{titlesec}
\usepackage{xcolor}
\usepackage{mdframed}
\usepackage[colorlinks=true,linkcolor=black,urlcolor=blue,
            citecolor=black]{hyperref}

%----------------------------- styling --------------------------------
\definecolor{glnavy}{RGB}{20,40,80}
\definecolor{glwarn}{RGB}{140,30,30}
\definecolor{glrule}{RGB}{200,205,215}

\titleformat{\section}{\Large\bfseries\color{glnavy}}{\thesection}{1em}{}
\titleformat{\subsection}{\large\bfseries}{\thesubsection}{1em}{}
\titleformat{\subsubsection}{\normalsize\bfseries}{\thesubsubsection}{1em}{}
\titlespacing*{\section}{0pt}{18pt}{8pt}

\pagestyle{fancy}
\fancyhf{}
\lhead{\small GuardianLink --- Video Run Sheet}
\rhead{\small\thepage}
\renewcommand{\headrulewidth}{0.4pt}

\setlist[itemize]{itemsep=2pt,topsep=4pt,parsep=0pt}
\setlist[enumerate]{itemsep=2pt,topsep=4pt,parsep=0pt}

\newcolumntype{L}[1]{>{\raggedright\arraybackslash}p{#1}}
\newcolumntype{C}[1]{>{\centering\arraybackslash}p{#1}}

% checkbox list
\newenvironment{checklist}
  {\begin{itemize}[label=$\square$,leftmargin=1.6em,itemsep=3pt]}
  {\end{itemize}}

% "say this" callout
\newcommand{\sayit}[1]{%
  \begin{mdframed}[linecolor=glrule,linewidth=0.8pt,
                   backgroundcolor=white,
                   innertopmargin=6pt,innerbottommargin=6pt,
                   innerleftmargin=8pt,innerrightmargin=8pt,
                   skipabove=6pt,skipbelow=6pt]
  \textit{Say:} ``#1''
  \end{mdframed}}

% warning callout
\newcommand{\warnbox}[1]{%
  \begin{mdframed}[linecolor=glwarn,linewidth=1pt,
                   innertopmargin=6pt,innerbottommargin=6pt,
                   innerleftmargin=8pt,innerrightmargin=8pt,
                   skipabove=8pt,skipbelow=8pt]
  \textbf{\color{glwarn}Warning.}\ #1
  \end{mdframed}}

%-----------------------------------------------------------------------
\begin{document}

\begin{center}
{\LARGE\bfseries\color{glnavy} GuardianLink\par}
\vspace{4pt}
{\Large Capstone Video --- Detailed Run Sheet\par}
\vspace{8pt}
{\small Target length: 12--15 minutes \quad$\cdot$\quad
 Student: [Your Name] \quad$\cdot$\quad Date: [Date]\par}
\end{center}
\vspace{6pt}
\hrule
\vspace{14pt}

\noindent
The brief asks the video to do one thing: demonstrate that every Must Have
works. The sequence below is ordered so each segment sets up the next ---
the forgot-password request is filed early so it is waiting in the admin
queue later, and every destructive action happens last so you cannot lock
yourself out mid-recording.

\tableofcontents
\vspace{10pt}
\hrule

%=======================================================================
\section{Before You Record}
%=======================================================================

\subsection{Start From a Clean Database}

A fresh seed gives predictable data and removes any half-broken accounts
left over from debugging.

\begin{Verbatim}[fontsize=\small,xleftmargin=1.5em]
cd "C:\Users\richa\WebstormProjects\HTML\GuardianLink Project"
node seed.js
node server.js
\end{Verbatim}

\subsection{Screen Hygiene}

\begin{checklist}
  \item Capture \textbf{one display only}. In OBS use Display Capture on the
        middle monitor and keep this run sheet open on another --- it will
        look like you are presenting from memory.
  \item \textbf{Close \texttt{.env} in WebStorm.} Never let
        \texttt{JWT\_SECRET} or an Atlas password appear on screen. The same
        applies to any terminal that has echoed them.
  \item Turn on Windows Focus Assist. Close Discord, Slack, and mail.
  \item New browser window: no extra tabs, bookmarks bar hidden, zoom at
        100\,\%.
  \item Open a \textbf{second browser profile or incognito window}. You will
        be logged in as two roles at once; this avoids constant
        log-out / log-in.
  \item Write the three demo passwords on paper, off-camera, so you do not
        fumble them live.
\end{checklist}

\subsection{Dry Run}

Click the entire sequence once with no narration. This surfaces the one
broken path before it is on tape, and settles the pacing.

\subsection{Recording Tool}

\begin{center}
\begin{tabularx}{\textwidth}{@{}L{3.2cm} X@{}}
\toprule
\textbf{Tool} & \textbf{Notes}\\
\midrule
OBS Studio   & Free, reliable, full control over which display is captured\\
Win+G        & Built into Windows, zero setup, adequate quality\\
Jumpshare    & Has a built-in recorder, which skips the upload step entirely\\
\bottomrule
\end{tabularx}
\end{center}

%=======================================================================
\section{Run Sheet Overview}
%=======================================================================

\begin{center}
\begin{tabularx}{\textwidth}{@{}C{0.8cm} L{5.2cm} C{1.6cm} X@{}}
\toprule
\textbf{\#} & \textbf{Segment} & \textbf{Time} & \textbf{Proves}\\
\midrule
0  & Intro                          & 0:30 & ---\\
1  & Landing / About                & 1:00 & About section requirement\\
2  & Locked out                     & 1:00 & Access control\\
3  & File forgot-password request   & 0:45 & Forgot password (part 1)\\
4  & Volunteer registration         & 1:30 & Volunteer data collection\\
5  & NGO registration               & 1:00 & NGO data collection\\
6  & NGO directory and contact      & 1:30 & Cross-role visibility, communication\\
7  & Volunteer directory            & 0:45 & Cross-role visibility (mirror)\\
8  & Self-service password change   & 1:00 & Extra credit --- security\\
9  & Admin panel                    & 3:00 & All admin privileges\\
10 & Self-deletion                  & 0:45 & Users delete own account only\\
11 & Mobile view                    & 0:45 & Responsive UI\\
12 & Code tour and GitHub           & 1:30 & Clean code, version control\\
13 & Close                          & 0:30 & ---\\
\midrule
\multicolumn{2}{@{}l}{\textbf{Total}} & \textbf{14:30} & \\
\bottomrule
\end{tabularx}
\end{center}

%=======================================================================
\section{Segment Detail}
%=======================================================================

\subsection*{0 --- Intro \hfill \normalfont\small 0:30}
\addcontentsline{toc}{subsection}{0 --- Intro}

Name, project name, one sentence on what it does, one sentence on the stack.

\sayit{GuardianLink connects volunteer cybersecurity specialists with
non-profits. It is Express and MongoDB serving a React frontend, with JWT
authentication and role-based access control.}

\subsection*{1 --- Landing Page \hfill \normalfont\small 1:00}
\addcontentsline{toc}{subsection}{1 --- Landing Page}

Scroll through and \textbf{name each required section out loud as you reach
it}. The grader is ticking boxes; make it effortless.

\sayit{Who we are. What we do. Why we do it. And our partners.}

\subsection*{2 --- Locked Out \hfill \normalfont\small 1:00}
\addcontentsline{toc}{subsection}{2 --- Locked Out}

Two parts, and the second is what separates this from most submissions.

\begin{enumerate}
  \item Paste the directory route while logged out --- you are redirected to
        login.
  \item Then open the API route directly in the address bar --- raw
        \texttt{401 Unauthorized} JSON.
\end{enumerate}

\begin{Verbatim}[fontsize=\small,xleftmargin=1.5em]
http://localhost:3000/#/directory      -> redirected to login
http://localhost:3000/api/directory    -> 401 Unauthorized
\end{Verbatim}

\sayit{The protection is not just the UI hiding things. Hitting the API
directly with no token returns 401 --- it is enforced server-side.}

\subsection*{3 --- Forgot Password, Part One \hfill \normalfont\small 0:45}
\addcontentsline{toc}{subsection}{3 --- Forgot Password, Part One}

Submit a reset request for the seeded NGO's email address. Mention that an
unknown email returns an identical confirmation, so the form cannot be used
to discover who is registered. Leave the request open --- you will resolve
it in segment 9.

\subsection*{4 --- Volunteer Registration \hfill \normalfont\small 1:30}
\addcontentsline{toc}{subsection}{4 --- Volunteer Registration}

Fill in every field the brief lists, naming each one as you type: name,
email, hours per week, background check, LinkedIn URL, areas of expertise.

Then \textbf{open MongoDB Compass and show the saved document}. Given the
\texttt{areas} / \texttt{hours} / \texttt{backgroundCheck} persistence bug
that was fixed during development, showing the array populated and the
boolean correct in the actual database is the strongest twenty seconds in
the video.

\subsection*{5 --- NGO Registration \hfill \normalfont\small 1:00}
\addcontentsline{toc}{subsection}{5 --- NGO Registration}

Point-of-contact email, organization name, areas of concern, password.

\subsection*{6 --- NGO Logged In \hfill \normalfont\small 1:30}
\addcontentsline{toc}{subsection}{6 --- NGO Logged In}

\begin{itemize}
  \item The directory shows volunteers and \textbf{no NGO cards} --- say so
        explicitly.
  \item Use the search box and the area filter.
  \item Open a volunteer profile, click \textbf{Contact}, and let the mail
        client open with recipient, subject, and body prefilled.
  \item Show the draft, then close it without sending.
\end{itemize}

\subsection*{7 --- Volunteer Logged In \hfill \normalfont\small 0:45}
\addcontentsline{toc}{subsection}{7 --- Volunteer Logged In}

Switch to the second browser window. NGOs only, no volunteers.

\sayit{Same rule, mirrored.}

\subsection*{8 --- Password Change \hfill \normalfont\small 1:00}
\addcontentsline{toc}{subsection}{8 --- Password Change}

\begin{enumerate}
  \item Enter a \textbf{wrong} current password --- rejected.
  \item Enter the correct one --- success, and you are logged out
        immediately.
  \item Log back in with the new password.
\end{enumerate}

Narrate it so the forced logout reads as deliberate rather than as a bug.

\sayit{Changing the password invalidates every token issued before it, so the
old session is dead on the next request.}

\subsection*{9 --- Admin Panel \hfill \normalfont\small 3:00}
\addcontentsline{toc}{subsection}{9 --- Admin Panel}

The longest segment. Follow this order exactly --- it is sequenced so you
never lock yourself out of the admin account.

\begin{enumerate}
  \item Log in as admin. The user list shows \textbf{all three roles},
        unlike the other two accounts.
  \item \textbf{Last-admin guard first}, while you are still the only admin.
        Try to demote yourself --- 409, nothing changes. Try to delete
        yourself --- the same.
  \item \textbf{Create a basic user profile.} Point out that your own
        session survives the operation.
  \item \textbf{Promote} the volunteer to admin. Note that the guard now
        relaxes, because two admins exist.
  \item \textbf{Demote} them back.
  \item \textbf{Resolve the reset request} filed in segment 3. Force-set a
        new password, and note that you were not asked for a current
        password, because an administrative override does not require one.
  \item \textbf{Delete} the throwaway user created in step 3.
\end{enumerate}

\warnbox{Do not improvise within this segment. Steps 2 and 7 are only safe
in these positions.}

\subsection*{10 --- Self-Deletion \hfill \normalfont\small 0:45}
\addcontentsline{toc}{subsection}{10 --- Self-Deletion}

Log in as a regular user. Point out that there is no Admin link in their
navigation and no delete control on anyone else's profile. Delete their own
account, then switch to the admin window and refresh to show the row gone.

\subsection*{11 --- Mobile View \hfill \normalfont\small 0:45}
\addcontentsline{toc}{subsection}{11 --- Mobile View}

DevTools, device toolbar, iPhone SE at 375\,px. Show the landing page, the
directory, and one form. No horizontal scrolling anywhere.

\subsection*{12 --- Code Tour and GitHub \hfill \normalfont\small 1:30}
\addcontentsline{toc}{subsection}{12 --- Code Tour and GitHub}

Four files, roughly twenty seconds each.

\begin{center}
\begin{tabularx}{\textwidth}{@{}L{4.6cm} X@{}}
\toprule
\textbf{File} & \textbf{Point out}\\
\midrule
\texttt{middleware/auth.js} & \texttt{requireAuth} and \texttt{requireRole}\\
\texttt{routes/directory.js} & The opposite-role filter derived from the
token, not from a query parameter\\
\texttt{models/User.js} & The \texttt{pre("save")} hash hook and its
\texttt{isModified} guard\\
\texttt{routes/admin.js} & The last-admin check\\
\bottomrule
\end{tabularx}
\end{center}

Then GitHub: the repository, the folder structure, and a scroll through the
commit history.

\subsection*{13 --- Close \hfill \normalfont\small 0:30}
\addcontentsline{toc}{subsection}{13 --- Close}

Thirty seconds on the extras: \texttt{helmet}, rate limiting,
\texttt{express-mongo-sanitize}, token invalidation on password change,
directory search and filtering, and the \texttt{seed.js} / \texttt{check.js}
tooling. Frame them as deliberate scope choices rather than scope creep.

%=======================================================================
\section{Pitfalls}
%=======================================================================

\begin{itemize}
  \item \textbf{Do not change any code during the recording.} If something
        breaks, stop, fix it, and re-record that segment.
  \item \textbf{Do not actually send the contact email.} Open the draft,
        show the prefilled fields, close it.
  \item \textbf{Say the status codes out loud} --- ``that is a 409
        Conflict'' --- whenever a guard fires. It shows the behaviour is
        designed rather than accidental.
  \item \textbf{Keep the server terminal in a corner of the capture} if the
        layout allows. Clean logs running alongside the UI are quietly
        convincing.
  \item \textbf{If you fluff a line, pause two seconds and say it again.}
        Far easier to cut than to re-record a whole segment.
\end{itemize}

%=======================================================================
\section{After Recording}
%=======================================================================

\begin{enumerate}
  \item Upload the MP4 to \url{https://jumpshare.com} --- free, and no
        account is needed for a basic share link.
  \item \textbf{Add timestamps to the description}, one per Must Have. A
        grader who can jump straight to ``Admin Privileges --- 8:15'' will
        find everything, which protects you if any single segment is weak.
  \item Paste the link into the capstone document title page, replacing the
        walk-through placeholder.
  \item \textbf{Confirm the link opens in a private browser window.} A link
        only you can view is the most common submission failure.
  \item Push a final commit so the repository matches what is in the video.
\end{enumerate}

%=======================================================================
\section*{Appendix --- Timestamp Template}
\addcontentsline{toc}{section}{Appendix --- Timestamp Template}
%=======================================================================

Fill these in after the final cut and paste them into the video description.

\begin{center}
\begin{tabularx}{\textwidth}{@{}L{2.2cm} X@{}}
\toprule
\textbf{Time} & \textbf{Must Have demonstrated}\\
\midrule
\rule{1.6cm}{0.4pt} & About page --- who, what, why, partners\\
\rule{1.6cm}{0.4pt} & Login required to view profiles\\
\rule{1.6cm}{0.4pt} & NGO registration and data capture\\
\rule{1.6cm}{0.4pt} & Volunteer registration and data capture\\
\rule{1.6cm}{0.4pt} & Cross-role visibility only\\
\rule{1.6cm}{0.4pt} & NGO-to-volunteer communication\\
\rule{1.6cm}{0.4pt} & Forgot password\\
\rule{1.6cm}{0.4pt} & Admin privileges\\
\rule{1.6cm}{0.4pt} & At least one admin at all times\\
\rule{1.6cm}{0.4pt} & Users delete own account only\\
\rule{1.6cm}{0.4pt} & Responsive desktop and mobile\\
\rule{1.6cm}{0.4pt} & Code walk-through and repository\\
\bottomrule
\end{tabularx}
\end{center}

\end{document}
